% concatenated from Grokability.tex
\documentclass[11pt]{article}

\usepackage[a4paper,margin=1in]{geometry}
\usepackage{amsmath,amssymb,amsthm,mathtools}
\usepackage{microtype}
\usepackage{hyperref}
\usepackage{xcolor}
\usepackage{enumitem}
\usepackage{verbatim}
\usepackage{authblk}

\newtheorem{theorem}{Theorem}
\newtheorem{remark}[theorem]{Remark}
\newtheorem{lemma}[theorem]{Lemma}

\hypersetup{
  colorlinks=true,
  linkcolor=blue!50!black,
  citecolor=blue!50!black,
  urlcolor=blue!50!black
}

% ---- macros & operators ----
\DeclareMathOperator{\sgn}{sgn}
\newcommand{\Maj}{\mathrm{Maj}}
\newcommand{\supp}
{\operatorname{supp}}

%\title{\bf  Mathematical Research Meets Large Language Models}
%\title{\bf Are We Close to an AI Mathematician?}
%\title{\bf Can Mathematical Research Be Accessible to Everyone?}
%\title{\bf Can Everyone Be a Mathematician in the AI Era?}
%\title{Reasoning Generative Models in Analytical Mathematics Research}
%\title{Reasoning Generative Systems in Analytical Mathematics Research}
%\title{Frontier Reasoning Models in Mathematical Research: New Results and Verification Challenges}
%\title{Frontier Reasoning LLMs in Mathematical Research}
%\title{Frontier Reasoning LLMs in Mathematical Research: New Results and Verification Challenges}
%\title{Frontier Reasoning LLMs in Mathematical Research: A Case Study with Grok}
\title{Grokability in five inequalities}


\author{
Paata Ivanisvili\thanks{University of California, Irvine, pivanisv@uci.edu} \quad  Xinyuan Xie\thanks{University of California, Irvine, xinyuax7@uci.edu} }


%\author{Ziang Chen\thanks{Massachusetts Institute of Technology, ziang@mit.edu} \quad
%Paata Ivanisvili\thanks{University of California, Irvine, pivanisv@uci.edu} \quad  Xinyuan Xie\thanks{University of California, Irvine, xinyuax7@uci.edu} }

\date{}

\begin{document}
\maketitle

\begin{abstract}
In this note, we report five mathematical discoveries made in collaboration with Grok, all of which have been subsequently verified by the authors. These include an improved lower bound on the maximal Gaussian perimeter of convex sets in $\mathbb{R}^n$, sharper $L_2$-$L_1$ moment comparison inequalities on the Hamming cube $\{-1,1\}^n$, a strengthened autoconvolution inequality, improved asymptotic bounds on the size of the largest $g$-Sidon sets in $\{1,\dots,n\}$, and an optimal balanced Szarek's inequality.
\end{abstract}


\section{Introduction}
In this note, we record several results in analysis, probability, convex geometry, and additive combinatorics, obtained through conversations with Grok. These results, to the best of our knowledge, have not previously appeared in the public literature. 


The first result is an improved lower bound for the maximal Gaussian perimeter of convex sets—the first constant improvement since Nazarov's construction in 2003 \cite{Nazarov2003}. The second result is an improved upper and lower bound on the optimal exponential base for $L_2-L_1$ moment comparison inequalities on the Hamming cube. In particular, the improved lower bound answers a MathOverflow question posed by Noam Lifshitz in 2014 \cite{MO184286}. The third result is an optimal Khintchine-type inequality when the Rademacher random variables are conditioned on the middle slice of the Hamming cube, and we call it balanced Szarek's inequality. In this case, the random variables are no longer independent, but interestingly, the classical semigroup approach can still obtain the optimal constant. The fourth result is an improved lower bound on an autocorrelation inequality related to the size estimate of $g$-Sidon sets. This constant is recorded as the $C_{1a}$ constant in a curated collection of optimization constants in mathematics initiated by Tao \cite{Tao26} and maintained by Davis, Ivanisvili, and Tao \cite{OptConstC1a}. This improvement is of interest from the prospective of  building automatic reasoning systems with the assistance of LLMs, as will be explained in the next paragraph. These examples, together with recent works obtained in a similar fashion \cite{IvanisviliXie25,BubeckCoesterEldanGowersLeeLupsascaSawhneyScherrerSellkeSpearsUnutmazWeilYinZhivotovskiy25,FengTrinhBinghamKangZhangKimBarretoSchildkrautJungSeoPaganoChervonyiHwangHouGukovTsaiChoiJinLiWuShiuShihLeLuong26, WoodruffCohenAddadJainMaoZuoBateniBranzeiBrennerChenFengFortnowFuGuanHadizadehHajiaghayiJafariRavizJavanmardKarthikKawarabayashiKumarLattanziLeeLiPanageasPaparasPrzybockiSubercaseauxSvenssonTaherijamWuYogevZadimoghaddamZhouMatiasManyikaMirrokni26,JuGaoJiangWuSunChenYWangYWangZWangHePWuXiaoLiuDaiDong26}, suggest that the current frontier reasoning models may carry out sophisticated analytical arguments, explore the boundary of existing methods, and implement new ideas with a meaningful degree of accuracy. We hope this note helps the community stay informed about the growing power of these new tools in mathematical exploration and discovery.

 
 On the other hand, there has been a series of works on building automated reasoning systems with the assistance of LLMs \cite{Georgiev25,Wang25,Yuksekgonul26}. For each problem, these systems first parameterize it, encode the parametrization in computer code, for example in Python, and then iteratively refine (evolve) the code with the help of LLMs. The idea of building automated reasoning systems is very interesting, but one may hope for a more ambitious use of LLMs. For example, the $C_{1a}$ constant has been extensively studied using computer-assisted approaches, whether based on LLMs or on classical implementations, for searches in explicit and structured discrete spaces; see \cite{OptConstC1a} for detailed information about this constant. In particular, the best known lower bound was obtained by exhausting certain large but finite sets using about 20{,}000 CPU hours \cite{CloningerSteinerberger17}. 
In this note, we present a tiny improvement to the lower bound for the $C_{1a}$ constant, obtained within a few minutes of conversation with Grok, suggesting that explicit and structural searching leaves extra room that can be uncovered through LLM-based conversational approaches\footnote{The term ``conversational approach'' was taken from a post by Tao \cite{Tao25Mathstodon}. }. This example, together with the other examples presented in this note, suggests the potential for building reasoning systems that explore mathematical questions through natural language and mathematical arguments. 


%Despite these successes, these models still make mathematical mistakes, especially in forms that are difficult for human users to detect. In particular, errors may arise when an argument involves large-scale numerics or hand-waving heuristics. It's important to illustrate the limitations of current models, and we present cautionary examples in which the model fails. If these models continue to make significant progress, then finding examples that are genuinely easy for humans to verify, but that the models still fail to solve, may become both valuable yet challenging.

The rest of the paper is organized as follows. In the remainder of the introduction, we present concise statements of each problem and explain why they are interesting. The proofs are presented in Section~\ref{proofs}. In Appendix~\ref{appendix:A}, we provide links that fully disclose our corresponding conversations with Grok, which led to the new results for each question.

 




%\begin{remark}[Disclosure of AI Use]
%The results reported here were first obtained through frontier reasoning models, and we provide links to fully disclose our corresponding conversations with Grok in Appendix \ref{appendix:A}. On the other hand, we found that most of these results can be recovered across many other thinking models, especially when an answer and a workable prompt are available. With these links, readers may find it interesting to challenge our prompts and obtain better results with your favorite models.
%\end{remark}

%a step-by-step logical progression: Looks naive. Pick author's overlook.



\subsection{Maximal Gaussian perimeter of convex sets}

For a convex set \(K\subseteq\mathbb{R}^n\), its Gaussian perimeter (or Gaussian surface area) is defined as
\[
\operatorname{GSA}(K)=\int_{\partial K}\varphi_n(x)\,d\sigma(x),
\qquad
\varphi_n(x):=(2\pi)^{-n/2}e^{-\|x\|^2/2}.
\]
The extremal quantity of interest is
\[
\Gamma(n):=\sup\{\operatorname{GSA}(K):K\subseteq\mathbb{R}^n \text{ is convex}\}.
\] The problem is to determine the growth of \(\Gamma(n)\) as \(n\to\infty\).

Ball\cite{Ball1993} proved that
\[
\frac{1}{e}\sqrt{\log n } \leq \Gamma(n)\le 4n^{1/4},
\]
Nazarov\cite{Nazarov2003} later showed that the order $n^{1/4}$ is sharp by proving
\[
e^{-5/4} < \liminf_{n\to\infty}\frac{\Gamma(n)}{n^{1/4}} \leq \limsup_{n\to\infty}\frac{\Gamma(n)}{n^{1/4}} < 0.64.
\]
The upper bound for $\Gamma_n$ was used to obtain sharper dimension dependence in multivariate Berry--Esseen bound over convex sets \cite{Bentkus03}. Rai\v{c} \cite{Raic2019} further refined the multivariate Berry--Essen bound and improved the upper bound of $\Gamma(n)$ to
\[
\Gamma(n)\le \sqrt{\frac{2}{\pi}}+0.59\bigl(n^{1/4}-1\bigr),
\] 
The problem is also relevant in learning theory. Klivans, O'Donnell, and Servedio showed that
Gaussian perimeter controls low-degree Hermite approximation, and hence the complexity of PAC
and agnostic learning under Gaussian distributions \cite{KOS2008}. More recently, Nadimpalli and Pascale\cite{NadimpalliPascale2025} revisited Nazarov's construction and recovered the Nazarov's lower-bound constant \(e^{-5/4}\) by an
alternative argument based on convex influence. Here we show that the constant in the lower bound can be improved further: 

\begin{theorem} \label{Gaussian} We have
    \[
\liminf_{n\to\infty}\frac{\Gamma_n}{n^{1/4}} \ \ge\ 0.31258.
\]
This is about $9\%$ improvement of the previous best known lower bound.  
\end{theorem}

In \cite{Nazarov2003}, the analytical argument is sophisticated, and several intermediate results are stated without proof. In Section~\ref{proofs}, we follow the same general strategy as in \cite{Nazarov2003}, while giving more detailed proofs and correcting a minor numerical inaccuracy in the original paper. We found that reasoning LLMs were able to work through these arguments, provide proofs for some statements left to the reader in \cite{Nazarov2003}, and fine-tune the parameters in Nazarov’s construction, leading to the present improvement.

\subsection{ $L_2$-$L_1$ moment comparison inequality}
For a function \(f:\{-1,1\}^n\to\mathbb{R}\) we write
\[
\|f\|_p:=\left(2^{-n}\sum_{x\in\{-1,1\}^n}|f(x)|^p\right)^{1/p},
\]
and we denote by \(\deg f\) its Fourier--Walsh degree, i.e.\ the degree of its multilinear expansion. We are interested in the optimal exponential base
\[
C_\ast:=\inf\Bigl\{C>0:\ \|f\|_2\le C^d\|f\|_1
\ \text{for every } d\ge 1,\ n\ge 1,\ \deg f\le d\Bigr\}.
\]

The problem of determining \(C_\ast\) was posed on MathOverflow by Noam Lifshitz in 2014, who asked in particular whether one can take \(C_\ast=\sqrt2\) \cite{MO184286}. The elementary example
\[
f(x)=\prod_{j=1}^d (1+x_j)
\]
already shows that \(C_\ast\ge \sqrt2\), since \(\|f\|_1=1\) and \(\|f\|_2=2^{d/2}\). The guess \(C_\ast=\sqrt2\) is also natural for structural reasons: for degree \(1\) it is exactly the sharp Khinchine inequality, due to Szarek \cite{Szarek76}; and if the same base \(\sqrt2\) were valid already for degree \(2\), then one would obtain the constant \(2\) for quadratic Rademacher chaoses, in accordance with the conjecture of Pe\l czy\'nski; see \cite[Remark~6]{IvanisviliTkocz}. On the upper-bound side, the real hypercontractivity yields \(C_\ast\le e\), see \cite[Theorem~9.22]{ODonnell}. A a more subtale argument via conformal maps, Hahn--Banach and Riesz representation theorem combined with complex hypercontractivity gives improvement $C_{*}\leq 2.69...$, see \cite{EskenazisIvanisvili}. We should also note that for \(d\)-homogeneous chaoses $h$ one has the sharper estimate
\[
\|h\|_2\le e^{d/2}\|h\|_1
\]
by Beckner's complex hypercontractivity \cite[Theorem~2]{IvanisviliTkocz}.

A useful observation is that if \(g\) is real-valued and \(f=g^2\), then
\[
\frac{\|f\|_2}{\|f\|_1}
=
\left(\frac{\|g\|_4}{\|g\|_2}\right)^2.
\]
Thus any family exhibiting large \(L^4/L^2\) growth immediately yields lower bounds for \(C_\ast\). The Gaussian model suggests taking \(g\) to be a Hermite polynomial evaluated on a long linear form. The next theorem shows that this already forces the base to be at least \(\sqrt3\).

\begin{theorem}\label{lem:sqrt3}
One has \(2.408...\geq C_\ast\ge \sqrt3\). 
\end{theorem}

We should note that the improvement in the upper bound from \(2.69\ldots\) to \(2.408\ldots\) was also known to the authors of \cite{EskenazisIvanisvili}. They did not revise the paper to reflect this, since the refinement follows verbatim from the argument in \cite{EskenazisIvanisvili}: one only needs to observe that the constant \(C\), defined in equation (191) there as the infimum of an explicit function, is in fact equal to \(2.408\ldots\).

On the lower-bound side, we note that the failure of the conjecture
\(C_\ast=\sqrt{2}\) in this generality was already known to some experts;
see, for instance, A. Samorodnitsky (private communication). Thus,
Theorem~\ref{lem:sqrt3} should not be regarded as a new mathematical fact
in itself. What is striking, however, is that reasoning models such as
Grok were able to recover both bounds, even though these observations do
not seem to have been recorded in the published literature.

\subsection{Optimal balanced Szarek's inequality}
The Khintchine inequality states that for independent Rademacher random variables
\(\varepsilon_1,\dots,\varepsilon_n\) and every \(p\geq 1\),
\[
A_p\Big(\sum_{i=1}^n a_i^2\Big)^{1/2}
\le
\Big(\mathbb{E}\Big|\sum_{i=1}^n a_i\varepsilon_i\Big|^p\Big)^{1/p}
\le
B_p\Big(\sum_{i=1}^n a_i^2\Big)^{1/2},
\]
where the constants \(A_p,B_p > 0\) depend only on \(p\). Determining the optimal constants in this inequality has attracted considerable attention: in
particular, the optimal constant when $p=1$ is due to Szarek \cite{Szarek76}, while the optimal constants for all $p$ were determined by Haagerup \cite{Haagerup81}.

A natural analogue is obtained by replacing the discrete cube \(\{-1,1\}^n\) with the middle
slice
\[
\Omega_n:=\Bigl\{x\in\{-1,1\}^n:\sum_{i=1}^n x_i=0\Bigr\},
\]
equipped with the uniform measure (here we need $n$ be even). Equivalently, one conditions the Rademacher random vector on the
event \(\sum_{i=1}^n x_i=0\) . In this model the coordinates are no longer independent, so the
classical proofs of Khintchine inequality may not directly apply. For \(p\ge 2\), such balanced Khintchine inequalities were studied by Spektor
\cite{Spektor16}. Later, Herscovici and Spektor \cite{HerscoviciSpektor20} computed the optimal constants when $p$ is even integer. By
contrast, the optimal constant for the endpoint \(p=1\) does not seem to have been documented in the literature. 

In the result below, we derive such balanced Khintchine-type iequality at the endpoint \(p=1\) with optimal constant, which may be viewed as a balanced analogue of Szarek's sharp \(L_1\) inequality.  Somewhat surprisingly, despite the loss of independence, the semigroup framework due to Kwapie\'n, Lata\l a, and Oleszkiewicz \cite{KwapienLatalaOleszkiewicz96,LatalaOleszkiewicz94,LatalaOleszkiewicz94} (see the second proof of Theorem 32 in \cite{NT23} for a modern presentation) still remains effective in obtaining the optimal constant. To deal with the lack of independence, we invoke a recent result of Filmus \cite{Filmus16} on orthogonal bases for the slice of the Hamming cube.

\begin{theorem}[Balanced Szarek's inequality]\label{Szarek}
Let \(n\in\mathbb{N}\) be even and \(n\ge 4\). Let \(x\) be distributed uniformly on
\[
\Omega_n:=\Bigl\{x\in\{-1,1\}^n:\sum_{i=1}^n x_i=0\Bigr\},
\]
and let \(a=(a_1,\dots,a_n)\in\mathbb{R}^n\). Then
\[
c_n \sqrt{\mathbb{E}\Big|\sum_{i=1}^n a_i x_i\Big|^2}
\le
\mathbb{E}\Big|\sum_{i=1}^n a_i x_i\Big|,
\]
where
\[
c_n=\sqrt{\frac{n-2}{2(n-1)}}
\]
is optimal.
\end{theorem}

\begin{remark}
The optimality is witnessed by the following examples, which is also the extremizer for classical Szarek's inequality. Let $a=(1, 1, 0,...,0)$. Then $|x_1+x_2|=2$ with probability $\frac{2\binom{n-2}{n/2 -2}}{\binom{n}{n/2}}=\frac{n-2}{2(n-1)}$, and $|x_1+x_2|=0$ otherwise. Thus, $\frac{\mathbb{E}|\sum_{i=1}^n a_i x_i|}{\sqrt{\mathbb{E}|\sum_{i=1}^n a_i x_i|^2}} = \sqrt{\frac{n-2}{2(n-1)}}$.
\end{remark}

\subsection{An autoconvolution inequality and $g$-Sidon sets}

Let an integrable $f : \mathbb{R} \to \mathbb{R}_{\ge 0}$ be supported on $[-1/4,1/4]$. We are interested in the largest universal constant $c>0$ such that
\[
\sup_{x\in\mathbb{R}} f*f(x)
\ge c \left( \int_{-1/4}^{1/4} f(x)\,dx \right)^2
\]
holds for every such function $f$. This constant \(c\) corresponds to the asymptotic size estimate for the  \(g\)-Sidon sets: a set \(A\subset\{1,2,\dots,n\}\) is called \(g\)-Sidon if every integer \(m\) has at most \(g\) representations of the form \(m=a+b\) with \(a,b\in A\). Denoting the largest size of a \(g\)-Sidon set for fixed $n$ and $g$ as
\[
\beta_g(n):=\max\{|A|:A\subset\{1,2,\dots,n\}\text{ is }g\text{-Sidon}\},
\]
\cite{MatolcsiVinuesa10} showed that the leading asymptotic constant for \(\beta_g(n)\) is a universal constant \(\sigma\) when $n$ and $g$ go to infinity, which can be reformulated in terms of the constant \(c\) as
\[
c=\frac{2}{\sigma^2}.
\]
Thus, studying the sharp autoconvolution inequality is equivalent to determining the leading asymptotic constant for the maximal size of \(g\)-Sidon sets. There have been many efforts devoted to bounding the constant \(c\) from above and below
\cite{SchinzelSchmidt02,MartinOBryant07,MartinOBryant09,MatolcsiVinuesa10,CloningerSteinerberger17,Georgiev25,Wang25,Yuksekgonul26}, where the best published bounds are
\[
1.28 \le c \le 1.50286.
\]

Recent progress on the upper bound \cite{Georgiev25,Wang25,Yuksekgonul26} is obtained through automatic reasoning systems assisted with LLMs, which iteratively modify Python codes describing candidate mathematical objects. At a high level, the reasoning LLMs in \cite{Georgiev25,Wang25,Yuksekgonul26} are used to explore the space formed by certain programming codes. One may naturally ask whether further improvements can still be achieved using existing techniques from the literature, and whether large language models can help uncover and exploit this remaining potential.  We provide positive evidence on this specific question.  In the following theorem, we provide an improved lower bound by refining an estimate in \cite{CloningerSteinerberger17}, which was suggested by LLMs.
\begin{theorem} \label{sidon}
 Let $f : \mathbb{R} \to \mathbb{R}_{\ge 0}$ be supported on $[-1/4,1/4]$. Then for every such function $f$, the following holds true
\[
\sup_{x\in\mathbb{R}} f*f(x)
\ge 1.2802 \left( \int_{-1/4}^{1/4} f(x)\,dx \right)^2.
\]   
\end{theorem}


The improvement can be obtained from Grok with a single natural prompt, simply by asking it to tighten an inequality in \cite{CloningerSteinerberger17}. The improvement is comparable to those obtained for the upper bound by LLM-based reasoning systems; see $C_{1a}$ constant \cite{OptConstC1a} for a neat record of recent progress on this question. From a mathematical perspective, however, this improvement is not particularly substantial, as it represents only a minor refinement of the estimate in \cite{CloningerSteinerberger17}. Moreover, we do not believe the authors missed anything essential, as their primary goal was to showcase the methodology. Nevertheless, we find this result impressive, since it shows that a purely conversational approach may extract a further improvement for a problem that has been extensively studied through structured and explicitly parameterized exhaustive search. 


%One may then wonder how far this dialogical approach can go. In the next section, we show that it can locate an improvement and implement the analytical argument in a more difficult situation, again within a couple of minutes.



%\subsection{Cautionary examples}
%Current advances still has limitations, suggests needs on better verification. 
%\begin{remark}
%With tools usage (being able to programming according to the reasoning process, amazing!), it becomes very challenging to find simple tasks
%\end{remark}
%1. an example related to NICD
%2. another example related to NICD


\section{Proofs}\label{proofs}



\subsection{Proof of Theorem~\ref{Gaussian}}
The proof presented in this section is the same as Nazarov's construction of a random convex set in \cite{Nazarov2003}. We show that a choice of specific parameters improves the previous lower bound. The random convex set is constructed in $\mathbb{R}^{n+1}$ as in Nazarov \cite{Nazarov2003} to avoid indexing $\varphi$, and we write $\Gamma_n:=\Gamma(n)$ for brevity.
\medskip



\begin{lemma}[Nazarov's exact expectation identity]\label{lem:ExpGSA}
Fix $n\ge 1$, $\rho>0$, and an integer $N\ge 1$.
Let $x_1,\dots,x_N$ be i.i.d.\ uniform on the unit sphere $S^n\subset\mathbb{R}^{n+1}$ and set
\[
Q:=\bigcap_{j=1}^N\{x\in\mathbb{R}^{n+1}:\ \langle x,x_j\rangle\le \rho\}.
\]
Then
\begin{equation}\label{eq:ExpGSA}
\mathbb{E}\,\operatorname{GSA}(Q)
=
N\,\frac{1}{\sqrt{2\pi}}e^{-\rho^2/2}
\int_{\mathbb{R}^n}\varphi_n(y)\,\bigl(1-p(|y|)\bigr)^{N-1}\,dy,
\end{equation}
where $p(r)$ depends only on $r\ge 0$ and is given by
\begin{equation}\label{eq:p_exact}
p(r)=
\Bigg(\int_{-\sqrt{r^2+\rho^2}}^{\sqrt{r^2+\rho^2}}
\Bigl(1-\frac{t^2}{r^2+\rho^2}\Bigr)^{\frac{n-2}{2}}\,dt\Bigg)^{-1}
\int_{\rho}^{\sqrt{r^2+\rho^2}}
\Bigl(1-\frac{t^2}{r^2+\rho^2}\Bigr)^{\frac{n-2}{2}}\,dt .
\end{equation}
\end{lemma}
\begin{remark}
In \cite{Nazarov2003}, the exponent of the integrand in the expression for $p(r)$ was given as $\frac{n-1}{2}$ instead of $\frac{n-2}{2}$ without proof, which is a minor inaccuracy that does not affect the final result. Here we include a proof for the reader's convenience.
\end{remark}
\begin{proof}
Since the laws of $x_j$'s are continuous and $x_j$ are independent, we have
$\mathbb{P}(x_i=x_j\text{ for some }i\neq j)=0$. Thus we may assume the normals
$x_1,\dots,x_N$ are pairwise distinct. For each $j$, let
\[
H_j:=\{x\in\mathbb{R}^{n+1}:\ \langle x,x_j\rangle=\rho\},
\qquad
F_j:=Q\cap H_j.
\]
Each $F_j$ is a (possibly empty) facet of $Q$ with outer unit normal $x_j$.
Moreover, $\partial Q$ is covered by the union of facets, and intersections
$F_i\cap F_j$ ($i\neq j$) have codimension at least $2$, hence $\sigma$-measure $0$.
Therefore,
\[
\operatorname{GSA}(Q)=\int_{\partial Q}\varphi_{n+1}(x)\,d\sigma(x)
=\sum_{j=1}^N\int_{F_j}\varphi_{n+1}(x)\,d\sigma(x).
\]
Taking expectation and using the fact that $\{x_i\}_{i \in [N]}$ has the same distribution,
\[
\mathbb{E}\,\operatorname{GSA}(Q)
=
N\cdot \mathbb{E}\!\left[\int_{F_1}\varphi_{n+1}(x)\,d\sigma(x)\right].
\]
By rotational invariance of $\varphi_{n+1}$, we may assume $x_1=e_{n+1}$. Therefore \[
\int_{F_1}\varphi_{n+1}(x)\,d\sigma(x) = \frac{1}{\sqrt{2\pi}}e^{-\rho^2/2}\int_{\mathbb{R}^n} \varphi_n(y) \mathbf{1}_{\{y:(y,\rho) \in F_1\}} dy,
\]
and by Tonelli, we have that 
\[
\mathbb{E}\!\left[\int_{F_1}\varphi_{n+1}(x)\,d\sigma(x)\right] =  \frac{1}{\sqrt{2\pi}}e^{-\rho^2/2}\int_{\mathbb{R}^n} \varphi_n(y) \mathbb{P}\big((y,\rho) \in F_1 \big) dy
\]
It remains to compute $\mathbb{P}\big((y,\rho) \in F_1 \big)$. For fixed $y\in\mathbb{R}^n$, the point $(y,\rho)$ belongs to $F_1$ iff it is not cut off
by any of the other half-spaces, i.e.
\[
\langle (y,\rho),x_j\rangle\le \rho,\qquad j=2,\dots,N.
\]
Let $r:=|y|$ and define
\[
p(r):=\mathbb{P}_{x\sim \mathrm{Unif}(S^n)}\big(\langle (y,\rho),x\rangle>\rho\big),
\]
which depends only on $|(y,\rho)|$ by rotation invariance of  $\mathrm{Unif}(S^n)$. Thus it only depends on $r$ for fixed $\rho$.
By independence,
\[
\mathbb{P}\big((y,\rho)\in F_1)=(1-p(r))^{N-1}.
\]
Hence
\[
\mathbb{E}\!\left[\int_{F_1}\varphi_{n+1}\,d\sigma\ \right]
=
\frac{1}{\sqrt{2\pi}}e^{-\rho^2/2}
\int_{\mathbb{R}^n}\varphi_n(y)\,(1-p(|y|))^{N-1}\,dy,
\]
and multiplying by $N$ gives \eqref{eq:ExpGSA}.

\emph{It remains to compute $p(r)$.} Fix $y$ with $|y|=r$ and put $u:=(y,\rho)\in\mathbb{R}^{n+1}$. For $x\sim \mathrm{Unif}(S^n)$, by rotational invariance, the random variable
$\frac{t}{|u|}:=\langle \frac{u}{|u|},x\rangle$ has the same distribution as $x_1$ as we may assume $\frac{u}{|u|}=e_1$. By \cite[Eq.~(8)]{BubeckDingEldanRacz16}, $x_1$ has density proportional to
\[
(1-|x_1|^2)^{(n-2)/2}.
\]
Therefore,
\[
p(r)=
\frac{\displaystyle\int_{\rho}^{|u|}
\left(1-\frac{t^2}{|u|^2}\right)^{(n-2)/2}\,dt}
{\displaystyle\int_{-|u|}^{|u|}
\left(1-\frac{t^2}{|u|^2}\right)^{(n-2)/2}\,dt}.
\]
Substituting $|u|=\sqrt{r^2+\rho^2}$ completes the proof. 
\end{proof}

\medskip

\begin{lemma}[Radialization]
 As in Nazarov \cite{Nazarov2003}, define
\begin{equation}\label{eq:f_c}
f(t):=t^{n-1}e^{-t^2/2},\qquad
c:=\Big(\int_0^\infty f(t)\,dt\Big)^{-1}.
\end{equation}
Then for every $n\geq 1$, every measurable $g:[0,\infty)\to[0,\infty]$,
\begin{equation}\label{eq:radialize}
\int_{\mathbb R^n}\varphi_n(y)\,g(|y|)\,dy
=
c\int_0^\infty f(r)\,g(r)\,dr.
\end{equation}
\end{lemma}
\begin{proof}
Let \(d\sigma\) denote the usual surface measure on \(S^{n-1}\). By polar coordinates and Tonelli,
\begin{align*}
\int_{\mathbb R^n}\varphi_n(y)g(|y|)\,dy
&=(2\pi)^{-n/2}\int_0^\infty\int_{S^{n-1}}
e^{-r^2/2}g(r)r^{n-1}\,d\sigma(\theta)\,dr \\
&=(2\pi)^{-n/2}|S^{n-1}|
\int_0^\infty r^{n-1}e^{-r^2/2}g(r)\,dr.
\end{align*}
Taking \(g\equiv 1\) gives
\[
1=(2\pi)^{-n/2}|S^{n-1}|\int_0^\infty f(r)\,dr,
\]
and hence
\[
(2\pi)^{-n/2}|S^{n-1}|=
\left(\int_0^\infty f(r)\,dr\right)^{-1}=c.
\]
Substituting this into the previous display gives the claim.
\end{proof}
\medskip


\begin{lemma}[Uniform upper bound for $p(r)$]\label{lem:p_bound}
Define the explicit Nazarov main term
\begin{equation}\label{eq:L_def}
L(n,\rho):=\frac{1}{\sqrt{2\pi}}\frac{1}{\rho}\exp\!\Big(\frac{\rho^4}{4n}\Big)e^{-\rho^2/2}.
\end{equation}
Fix constants $a>0$ and $W>0$. Then there exist constants $C_{a,W}>0$ and
$n_0(a,W)$ such that for all $n\ge n_0(a,W)$, all $\rho=a n^{1/4}$, and all
$r=\sqrt{n-1}+w$ with $|w|\le W$, one has
\begin{equation}\label{eq:p_bound}
p(r)\ \le\ L(n,\rho)\,\exp\!\Big(\frac{w\rho^2}{\sqrt{n}}\Big)\,
\exp\!\Big(\frac{C_{a,W}}{\sqrt{n}}\Big).
\end{equation}
\end{lemma}

\begin{proof}
Put $S:=r^2+\rho^2$. Then the denominator in \eqref{eq:p_exact} equals
\[
\sqrt S\int_{-1}^{1}(1-u^2)^{\frac{n-2}{2}}du
=
\sqrt S\,\sqrt\pi\,\frac{\Gamma(n/2)}{\Gamma((n+1)/2)}.
\]
By Wendel's inequality~\cite{Wendel1948},
\[
\frac{\Gamma(n/2)}{\Gamma((n+1)/2)}
\ge \sqrt{\frac{2}{n}},
\]
and hence
\[
\sqrt S\,\sqrt\pi\,\frac{\Gamma(n/2)}{\Gamma((n+1)/2)}
\ge \sqrt{\frac{2\pi S}{n}} .
\]
For the numerator, use
\[
\log(1-u)\le -u-\frac{u^2}{2},\qquad 0<u<1.
\]
Thus, for \(0<t<\sqrt S\),
\[
\Bigl(1-\frac{t^2}{S}\Bigr)^{\frac{n-2}{2}}
\le
\exp\!\left(
-\frac{n-2}{2S}t^2-\frac{n-2}{4S^2}t^4
\right).
\]
Since \(t\ge \rho\) on the interval of integration,
\[
\int_{\rho}^{\sqrt S}
\Bigl(1-\frac{t^2}{S}\Bigr)^{\frac{n-2}{2}}dt
\le
\exp\!\left(-\frac{n-2}{4S^2}\rho^4\right)
\int_{\rho}^{\infty}
\exp\!\left(-\frac{n-2}{2S}t^2\right)dt .
\]
Set 
\[
m(x):=\frac{1-\Phi(x)}{\varphi_1(x)}
=
\frac{\int_x^\infty e^{-u^2/2}\,du}{e^{-x^2/2}},
\qquad x>0.
\]
It follows from integration by parts that  \(m(x)\le x^{-1}\), i.e.,  equivalently,
\[
\int_x^\infty e^{-u^2/2}\,du\le \frac{1}{x}e^{-x^2/2}.
\qquad x>0,
\]
 Applying this with
\[
x=\rho\sqrt{\frac{n-2}{S}},
\]
we get, for \(n\ge3\),
\[
\int_{\rho}^{\infty}
\exp\!\left(-\frac{n-2}{2S}t^2\right)dt
=
\sqrt{\frac{S}{n-2}}
\int_{\rho\sqrt{(n-2)/S}}^\infty e^{-u^2/2}\,du
\le
\frac{S}{(n-2)\rho}
\exp\!\left(-\frac{n-2}{2S}\rho^2\right).
\]
Therefore
\[
\int_{\rho}^{\sqrt S}
\Bigl(1-\frac{t^2}{S}\Bigr)^{\frac{n-2}{2}}dt
\le
\frac{S}{(n-2)\rho}
\exp\!\Big(
-\frac{n-2}{2S}\rho^2-\frac{n-2}{4S^2}\rho^4
\Big).
\] Hence
\[
p(r)\le
\frac{1}{\sqrt{2\pi}\rho}\,
\frac{S}{n-2}\sqrt{\frac{n}{S}}\,
\exp\!\Big(
-\frac{n-2}{2S}\rho^2-\frac{n-2}{4S^2}\rho^4
\Big).
\]
Now set $\rho=a n^{1/4}$ and $r=\sqrt{n-1}+w$, with $|w|\le W$. Then
\[
S=r^2+\rho^2
=(\sqrt{n-1}+w)^2+\rho^2
=n+(2w+a^2)\sqrt n+O_{a,W}(1),
\]
uniformly for $|w|\le W$. Hence
\[
\frac{n-2}{S}
=
1-\frac{2w+a^2}{\sqrt n}+O_{a,W}(n^{-1}),
\qquad
\frac{n-2}{S^2}
=
\frac1n+O_{a,W}(n^{-3/2}).
\]
Since $\rho^2=a^2\sqrt n$ and $\rho^4=a^4 n$, we get
\[
-\frac{n-2}{2S}\rho^2
=
-\frac{\rho^2}{2}
+\frac{w\rho^2}{\sqrt n}
+\frac{\rho^4}{2n}
+O_{a,W}(n^{-1/2}),
\]
and
\[
-\frac{n-2}{4S^2}\rho^4
=
-\frac{\rho^4}{4n}
+O_{a,W}(n^{-1/2}).
\]
Adding the last two displays gives
\[
-\frac{n-2}{2S}\rho^2-\frac{n-2}{4S^2}\rho^4
=
-\frac{\rho^2}{2}
+\frac{w\rho^2}{\sqrt n}
+\frac{\rho^4}{4n}
+O_{a,W}(n^{-1/2}),
\]
uniformly for $|w|\le W$. Substituting the exponent expansion together with
\[
\frac{S}{n-2}\sqrt{\frac{n}{S}}
=
\exp\bigl(O_{a,W}(n^{-1/2})\bigr)
\]
gives
\[
p(r)\le
\frac{1}{\sqrt{2\pi}}\frac1\rho
\exp\!\Big(-\frac{\rho^2}{2}+\frac{\rho^4}{4n}\Big)
\exp\!\Big(\frac{w\rho^2}{\sqrt n}\Big)
\exp\!\Big(\frac{C_{a,W}}{\sqrt n}\Big),
\]
which is exactly \eqref{eq:p_bound} by the definition of $L(n,\rho)$. Increasing
$n_0(a,W)$ if necessary so that $n\ge3$ completes the proof.
\end{proof}



\medskip

\begin{lemma}[Uniform Laplace form for $c\,f(\sqrt{n-1}+w)$]\label{lem:cf}
Fix $W>0$. There exist constants $C_W>0$ and $n_1(W)$ such that for all
$n\ge n_1(W)$ and all $|w|\le W$,
\begin{equation}\label{eq:cf_asymp}
c\,f(\sqrt{n-1}+w)
=
\frac{1}{\sqrt{\pi}}e^{-w^2}\exp\!\Big(\frac{\theta_{n,W}(w)}{\sqrt{n}}\Big),
\qquad |\theta_{n,W}(w)|\le C_W.
\end{equation}
\end{lemma}

\begin{proof}
Let $t_0:=\sqrt{n-1}$ and $u:=w/t_0$. 
Then
\[
\log\frac{f(t_0+w)}{f(t_0)}=(n-1)\log(1+u)-t_0w-\frac{w^2}{2}.
\]
For $n\ge 4W^2+1$ we have $|u|\le 1/2$. Using $|\log(1+u)-(u-u^2/2)|\le 2|u|^3$ for $|u|\le 1/2$ and the identities
$(n-1)u=t_0w$, $(n-1)u^2=w^2$, we obtain
\[
\log\frac{f(t_0+w)}{f(t_0)}=-w^2+O_W(n^{-1/2})
\quad\text{uniformly for }|w|\le W.
\]
Next, $\int_0^\infty f(t)\,dt=2^{n/2-1}\Gamma(n/2)$, so $c f(t_0)$ can be
estimated by Stirling bounds for $\Gamma(n/2)$, giving
$c f(t_0)=\frac{1}{\sqrt{\pi}}\exp(O(n^{-1}))$.
Combining the two estimates yields \eqref{eq:cf_asymp}.
\end{proof}

\medskip

\begin{proof}[Proof of Theorem~\ref{Gaussian}]
Fix the rational parameters
\[
a:=\frac{6131}{5000}=1.2262,
\qquad
b:=\frac{2387}{1000}=2.387,
\qquad
W:=6.
\]
For each $n$, set
\[
\rho:=a n^{1/4},
\qquad
N:=\Big\lfloor \frac{b}{L(n,\rho)}\Big\rfloor,
\]
where $L(n,\rho)$ is defined in \eqref{eq:L_def}. Since $L(n,\rho)\to 0$
exponentially fast in $\sqrt{n}$, we have
\begin{equation}\label{eq:NL_to_b}
N\,L(n,\rho)\to b
\qquad\text{as }n\to\infty.
\end{equation}

Let $Q\subset\mathbb{R}^{n+1}$ be the random polytope from Lemma~\ref{lem:ExpGSA}.
By definition of $\Gamma(n+1)$,
\[
\Gamma(n+1)\ \ge\ \mathbb{E}\,\operatorname{GSA}(Q).
\]
Using Lemma~\ref{lem:ExpGSA} and radialization \eqref{eq:radialize},
\begin{equation}\label{eq:ExpGSA_radial}
\mathbb{E}\,\operatorname{GSA}(Q)
=
N\,\frac{1}{\sqrt{2\pi}}e^{-\rho^2/2}\,
c\int_0^\infty f(r)\bigl(1-p(r)\bigr)^{N-1}\,dr.
\end{equation}
Restrict to $r=\sqrt{n-1}+w$ with $|w|\le W$:
\begin{equation}\label{eq:restrictW}
\mathbb{E}\,\operatorname{GSA}(Q)
\ge
N\,\frac{1}{\sqrt{2\pi}}e^{-\rho^2/2}
\int_{-W}^{W} c\,f(\sqrt{n-1}+w)\,
\bigl(1-p(\sqrt{n-1}+w)\bigr)^{N-1}\,dw.
\end{equation}

By Lemma~\ref{lem:p_bound}, for $|w|\le W$ and large $n$,
\[
p(\sqrt{n-1}+w)\le q_n(w):=
L(n,\rho)\exp\!\Big(\frac{w\rho^2}{\sqrt{n}}\Big)\exp\!\Big(\frac{C_{a,W}}{\sqrt{n}}\Big).
\]
Since $q_n(w)\to 0$ uniformly on $[-W,W]$, for large $n$ we have $q_n(w)\le 1/2$ there,
and $\log(1-x)\ge -x-x^2$ on $[0,1/2]$ gives
\[
(1-q_n(w))^{N-1}\ge \exp\big(-(N-1)q_n(w)-(N-1)q_n(w)^2\big).
\]
Using $\rho^2/\sqrt{n}=a^2$ and \eqref{eq:NL_to_b}, we get pointwise for each fixed $w$,
\[
(N-1)q_n(w)\to b e^{a^2 w},
\qquad
(N-1)q_n(w)^2\to 0,
\]
hence
\begin{equation}\label{eq:lim_factor2}
\liminf_{n\to\infty}\bigl(1-p(\sqrt{n-1}+w)\bigr)^{N-1}
\ \ge\ \exp\big(-b e^{a^2 w}\big).
\end{equation}
Also Lemma~\ref{lem:cf} gives
\begin{equation}\label{eq:lim_factor1}
c\,f(\sqrt{n-1}+w)\to \frac{1}{\sqrt{\pi}}e^{-w^2} \qquad\text{as }n\to\infty.
\end{equation}

Next, by the definition \eqref{eq:L_def} we have the exact identity
\[
\frac{1}{\sqrt{2\pi}}e^{-\rho^2/2}=\rho e^{-\rho^4/(4n)}\,L(n,\rho).
\]
Therefore, using \eqref{eq:NL_to_b} and $\rho=a n^{1/4}$,
\begin{equation}\label{eq:prefactor_limit}
\frac{N}{n^{1/4}}\frac{1}{\sqrt{2\pi}}e^{-\rho^2/2}
=
\frac{N L(n,\rho)}{n^{1/4}}\rho e^{-\rho^4/(4n)}
\to
b\,a\,e^{-a^4/4} \qquad\text{as }n\to\infty.
\end{equation}

Apply Fatou's lemma to \eqref{eq:restrictW} and use
\eqref{eq:lim_factor2}, \eqref{eq:lim_factor1}, \eqref{eq:prefactor_limit}. We obtain
\begin{equation}\label{eq:liminf_W}
\liminf_{n\to\infty}\frac{\Gamma(n+1)}{n^{1/4}}
\ge
b\,a\,e^{-a^4/4}\cdot \frac{1}{\sqrt{\pi}}
\int_{-W}^{W}\exp\big(-w^2-b e^{a^2 w}\big)\,dw.
\end{equation}

\smallskip\noindent
\emph{Certified numerical lower bound (with $W=6$).}
Let
\[
F(w):=\exp\big(-w^2-b e^{a^2 w}\big),
\qquad
J:=\int_{-6}^{6}F(w)\,dw.
\]
We lower bound $J$ by the composite trapezoidal rule with step $h:=5\cdot 10^{-4}$.
Let $N_h:=24000$ and $w_k:=-6+kh$ ($0\le k\le N_h$). Define
\[
T:=h\Big(\tfrac12 F(w_0)+\sum_{k=1}^{N_h-1}F(w_k)+\tfrac12 F(w_{N_h})\Big).
\]
By the standard error estimate for the composite trapezoidal rule
\cite[Sec.~4.4, Thm.~4.5]{BurdenFairesBurden2016}, if
\(f\in C^2([A,B])\), \(h=(B-A)/m\), and
\[
T_h
=
h\left(\frac{f(A)+f(B)}{2}
+\sum_{k=1}^{m-1} f(A+kh)\right),
\]
then there exists \(\xi\in(A,B)\) such that
\[
\int_A^B f(x)\,dx-T_h
=
-\frac{B-A}{12}h^2 f''(\xi).
\]
In particular,
\[
\left|\int_A^B f(x)\,dx-T_h\right|
\le
\frac{B-A}{12}h^2\|f''\|_{L^\infty[A,B]}.
\]
Applying this with \(A=-6\), \(B=6\), \(f=F\), and \(m=24000\), gives
\[
|J-T|\le \frac{12}{12}h^2\max_{[-6,6]}|F''|
=h^2\max_{[-6,6]}|F''|.
\]
Writing $F=e^{g}$ with $g(w)=-w^2-b e^{a^2 w}$, one has
$F''=(g''+(g')^2)F$. A termwise global estimate using
$\sup_{t>0}t e^{-bt}=1/(be)$,
$\sup_{t>0}t^2 e^{-bt}=(2/b)^2e^{-2}$,
$\sup_{w\in\mathbb{R}}w^2e^{-w^2}=1/e$,
$\sup_{w\in\mathbb{R}}|w|e^{-w^2}=1/\sqrt{2e}$
yields, for all $w\in\mathbb{R}$,
\[
|F''(w)|\le
M(a):=
2+\frac{a^4}{e}+\frac{4}{e}+\frac{4a^2}{e\sqrt{2e}}+4a^4 e^{-2}.
\]
For $a=6131/5000$ one checks $M(a)<6.476$.

A direct finite computation gives
\[
T \geq 0.33310717054594
\]
and therefore
\[
J\ge T - M(a)h^2 \ \ge\ 0.33310555154594.
\]
Moreover,
\[
b\,a\,e^{-a^4/4} \geq 1.6632596039 \qquad
\frac{1}{\sqrt{\pi}}\ \ge\ 0.5641895835.
\]
Hence the right-hand side of \eqref{eq:liminf_W} with $W=6$ is at least
\[
1.6632596039 \times 0.5641895835 \times 0.33310555154594 >\ 0.312584.
\]
Finally, since $(n/(n+1))^{1/4}\to 1$, this implies
\[
\liminf_{n\to\infty}\frac{\Gamma(n)}{n^{1/4}}\ \ge\ 0.312584,
\]
which is exactly Theorem~\ref{Gaussian}.
\end{proof}


\subsection{Proof of Theorem~\ref{lem:sqrt3}}

\begin{proof}[Proof of Theorem~\ref{lem:sqrt3}]
Let \(\gamma\) be the standard Gaussian measure on \(\mathbb{R}\), and let \(\mathrm{He}_m\) denote the monic probabilists' Hermite polynomial of degree \(m\), normalized so that
\[
\int_{\mathbb{R}} \mathrm{He}_m(x)\mathrm{He}_k(x)\,d\gamma(x)=m!\,\delta_{mk}.
\]
In particular,
\[
\|\mathrm{He}_m\|_{L^2(\gamma)}^2=m!.
\]

For \(N\ge 1\), define
\[
S_N(x):=\frac{x_1+\cdots+x_N}{\sqrt N},
\qquad x\in\{-1,1\}^N,
\]
and
\[
F_{m,N}(x):=\mathrm{He}_m(S_N(x))^2.
\]
Since \(\mathrm{He}_m\) has degree \(m\), the function \(\mathrm{He}_m(S_N(x))\) is a polynomial of total degree at most \(m\) in the coordinates \(x_1,\dots,x_N\). After multilinear reduction on the cube (using \(x_i^2=1\)), its Fourier--Walsh degree is still at most \(m\). Hence
\[
\deg F_{m,N}\le 2m.
\]

If \(\varepsilon_1,\dots,\varepsilon_N\) are i.i.d.\ Rademacher signs, then
\[
\|F_{m,N}\|_1=\mathbb{E}\,\mathrm{He}_m\!\left(\frac{\varepsilon_1+\cdots+\varepsilon_N}{\sqrt N}\right)^2,
\]
and
\[
\|F_{m,N}\|_2
=
\left(
\mathbb{E}\,\mathrm{He}_m\!\left(\frac{\varepsilon_1+\cdots+\varepsilon_N}{\sqrt N}\right)^4
\right)^{1/2}.
\]
Now
\[
\frac{\varepsilon_1+\cdots+\varepsilon_N}{\sqrt N}\ \Longrightarrow\ G
\qquad(N\to\infty),
\]
where \(G\sim N(0,1)\), and in fact all moments converge. Since \(\mathrm{He}_m^2\) and \(\mathrm{He}_m^4\) are fixed polynomials, it follows that
\[
\lim_{N\to\infty}\frac{\|F_{m,N}\|_2}{\|F_{m,N}\|_1}
=
\frac{\|\mathrm{He}_m\|_{L^4(\gamma)}^2}{\|\mathrm{He}_m\|_{L^2(\gamma)}^2}.
\]

We now use the asymptotics of Larsson-Cohn. By \cite[Theorem~2.1(b)]{LarssonCohn}, for every fixed \(p>2\),
\[
\|\mathrm{He}_m\|_{L^p(\gamma)}
=
c(p)\,m^{-1/4}\sqrt{m!}\,(p-1)^{m/2}(1+o(1))
\qquad(m\to\infty),
\]
where \(c(p)>0\) is an explicit constant. Taking \(p=4\), we obtain
\[
\|\mathrm{He}_m\|_{L^4(\gamma)}
=
c(4)\,m^{-1/4}\sqrt{m!}\,3^{m/2}(1+o(1)),
\]
and therefore
\[
\frac{\|\mathrm{He}_m\|_{L^4(\gamma)}^2}{\|\mathrm{He}_m\|_{L^2(\gamma)}^2}
=
c(4)^2\,m^{-1/2}\,3^m(1+o(1)).
\]
Consequently,
\[
\lim_{m\to\infty}
\left(
\frac{\|\mathrm{He}_m\|_{L^4(\gamma)}^2}{\|\mathrm{He}_m\|_{L^2(\gamma)}^2}
\right)^{1/(2m)}
=
\sqrt3.
\]

Let \(C<\sqrt3\). By the previous limit, one can choose \(m\) so large that
\[
\frac{\|\mathrm{He}_m\|_{L^4(\gamma)}^2}{\|\mathrm{He}_m\|_{L^2(\gamma)}^2}>C^{2m}.
\]
Fix such an \(m\). Then, by the convergence in \(N\), for all sufficiently large \(N\),
\[
\frac{\|F_{m,N}\|_2}{\|F_{m,N}\|_1}>C^{2m}.
\]
But \(\deg F_{m,N}\le 2m\), so the inequality
\[
\|f\|_2\le C^d\|f\|_1
\]
fails for \(d=2m\). Since this happens for every \(C<\sqrt3\), we conclude that \(C_\ast\ge \sqrt3\).
\end{proof}

\subsection{Proof of Theorem~\ref{Szarek}}
Let $n \in \mathbb{N}$ be even. In this section, we show that the balanced Szarek's inequality, i.e.,  $(x_1,...,x_n)$ is uniformly distributed on $\Omega_n:=\{(x_1,...,x_n) \in \{-1,1\}^n:\sum_{i=1}^n x_i =0\}$, holds with the optimal constant.

The derivation of the optimal constant follows the semigroup proof framework for Szarek's inequality (also called the \(L_1\) Khintchine inequality) due to Kwapie\'n, Lata\l a, and Oleszkiewicz \cite{KwapienLatalaOleszkiewicz96,LatalaOleszkiewicz94,LatalaOleszkiewicz94} (see the second proof of Theorem 32 in \cite{NT23} for a modern presentation). In their setting, the random variables are independent, whereas in our setting they are slightly dependent. Their framework, however, still works. To deal with the lack of independence, we invoke a recent result of Filmus \cite{Filmus16} on orthogonal bases for the slice of the Hamming cube.

This may suggest that a certain symmetry of the distribution can be a suitable replacement for the independence condition in some Khintchine-type inequalities. 

\subsubsection{Harmonic Multilinear Polynomial}
Let $f:\mathbb{R}^n \to \mathbb{R}$ be a multilinear polynomial. We call $f$ a harmonic multilinear polynomial if 
\[
\sum_{i=1}^n \frac{\partial f}{\partial x_i}(x)=0 \quad \text{for all } x \in \mathbb{R}^n.
\]
One can check that the harmonic multilinear polynomials on $\mathbb{R}^n$ form a vector space, which we denote by $H_n$. Note that $H_n$ has dimension $\binom{n}{\lfloor\frac{n}{2}\rfloor}$ and every element in $H_n$ is a multilinear polynomial of degree at most $\frac{n}{2}$; see \cite[Lemma 2]{Filmus16} for a proof. 

A probability distribution $\mu$ on $\mathbb{R}^n$ is called exchangeable if it is invariant under permutations of the coordinates. Given an exchangeable distribution $\mu$ with mild regular conditions, we define the inner product on $H_n$ by 
\[
\langle f,g \rangle=\mathbb{E}_{x\sim\mu}[f(x)g(x)], 
\]
and set $\| f \|_2^2 = \langle f,f \rangle$. 
Filmus \cite[Section 3]{Filmus16} constructed a finite collection of functions $\{\chi_B\}_{B \in B_n}$ that forms an orthogonal basis of $H_n$ with respect to any exchangeable probability distribution $\mu$, and showed that every $f\in H_n$ admits the following Young-Fourier expansion  (whenever $\|\chi_B\|_2^2 > 0$ for all $B \in B_n$):
\[
f = \sum_{B \in B_n}\widehat f(B)\chi_B,
\]
where $\widehat f(B)=\frac{\langle f,\chi_B \rangle}{\|\chi_B\|_2^2}$. We will use the following facts from \cite[Section 3]{Filmus16}: every $\chi_B$ in the basis is a homogeneous multilinear polynomial indexed by $B$, where $B$ is a list of increasing numbers in $[n]$; the length of $B$, denoted by $|B|$, equals the degree of $\chi_B$.  

Filmus \cite[Section 6]{Filmus16} further defined a Laplacian operator $L$ by 
\[
Lf = f - \frac{1}{\binom{n}{2}}\sum_{1\leq i<j \leq n}f^{(i,j)},
\]
where $f^{(i,j)}=f(x^{(i,j)})$ and $x^{(i,j)}$ is obtained from $x$ by swapping its $i$th and $j$th coordinates. He also showed that $L$ admits the following spectral decomposition:
\[
Lf = \sum_{B \in B_n} \frac{2|B|(n+1-|B|)}{n(n-1)}\widehat f(B) \chi_B.
\]


\subsubsection{Balanced Szarek's inequality with best constant}
\begin{theorem}\label{uniqueness}
Let $n \in \mathbb{N}$ be even. For any function $f:\Omega_n \to \mathbb{R}$, there exists a unique $\tilde f \in H_n$ such that $\tilde f(x)= f(x)$ for all $x \in \Omega_n$.
\end{theorem}
\begin{proof}
Since the dimension of functions on $\Omega_n$ is $\binom{n}{\frac{n}{2}}$, which matches the cardinality of $\{\chi_B\}_{B\in B_n}$, it suffices to show that $\{\chi_B\}_{B\in B_n}$ is a set of linearly independent vectors when restricted to $\Omega_n$. Equip $\mathbb{R}^n$ with the uniform probability measure on $\Omega_n$. Thus it suffices to show that $\{\chi_B\}_{B\in B_n}$ forms an orthogonal basis with respect to the uniform probability measure on $\Omega_n$. Since the uniform probability measure on $\Omega_n$ is exchangeable, by Theorem 9 in \cite{Filmus16}, we know that $\langle \chi_A,\chi_B \rangle =0 $ for any $A,B \in B_n$ with $A \neq B$. It remains to show that $\| \chi_B\|^2_2> 0$ for all $B \in B_n$. Theorem 10 in \cite{Filmus16} shows that, for $B \in B_n$ of length $d$,
\[
\| \chi_B\|_2^2=c_B\|\chi_d\|_2^2
\]
where $\chi_d=\prod_{i=1}^d (x_{2i-1}-x_{2i})$ and $c_B= \prod_{i=1}^d \frac{(b_i - 2(i-1))(b_i - 2(i-1)-1)}{2} $. By the definition of $B \in B_n$ (see the definition of the {\em top} set under the Definition 4 in \cite{Filmus16}) , one can deduce that $b_i \geq 2i$ for all $i \in [d]$. Hence $c_B>0$ and it suffices to show that $\| \chi_d\|_2^2 >0.$ Let $a\in \Omega_n$ such that $a_i=(-1)^i$ for all $i \in [n]$. Since $\chi_d(a) \neq 0$ and $\Pr\{x =a\}>0$ under the uniform probability measure on $\Omega_n$, we conclude that $\|\chi_d\|_2^2 > 0$, completing the proof. 
\end{proof}
For the rest of the section, we equip $\mathbb{R}^n$ with the uniform probability measure on $\Omega_n$, which is exchangeable. Hence, many of the results in \cite{Filmus16} apply. Let $n \in \mathbb{N}$ be even, and let $(a_1,...,a_n) \in \mathbb{R}^n$. By Theorem~\ref{uniqueness}, there exists $f \in H_n$ such that
\[
f(x)=|\sum_{i=1}^n a_i x_i| \quad \text{for all } x \in \Omega_n.
\]
\begin{theorem}[Spectrum is supported on even indices] \label{eveness}
The spectrum of $f$ defined above is supported on the even indices, i.e., 
\[
f(x) = \sum_{\substack{B \in B_n\\ |B| \text{ even}}} \widehat{f}(B)\chi_B(x).
\]    
\end{theorem}
\begin{proof}
Note that $f(x)=f(-x)$ for all $x \in \Omega_n$. Looking at this fact from the Young-Fourier expansion of $f$, we have that
\[
\sum_{B \in B_n} \widehat f(B)\chi_B(x)= \sum_{B \in B_n}\widehat{f}(B)(-1)^{|B|}\chi_B(x), \quad \text{for all } x \in \Omega_n
\]
due to the fact that $\chi_B$ is a homogeneous multilinear polynomial of degree $|B|$. After rearranging the terms, we have that 
\[
\sum_{\substack{B \in B_n\\|B| \text{ is odd}}} \widehat{f}(B)\chi_B(x)=0,\quad \text{for all }x \in \Omega_n.
\]
Since the functions $\{\chi_B\}_{B\in B_n}$ are linearly independent by Theorem~\ref{uniqueness}, we know that $\widehat{f}(B)=0$ for all $B \in B_n$ with $|B|$ odd, completing the proof.     
\end{proof}


A fact used in the Kwapie\'n-Lata\l a-Oleszkiewicz's proof of the Szarek's inequality is a tightened Poincar\'e inequality for $f$ discussed above. We now derive an analogue on the middle slice. 
\begin{theorem}[A tightened Poincar\'e inequality] \label{Poincare}
 Let $n\in \mathbb{N}$ be even. Let $f\in H_n$ such that $f(x)=|\sum_{i=1}^n a_ix_i|$ for all $x \in \Omega_n$.  $\mathrm{Var}(f):=\mathbb{E}_{\mathrm{unif}\{\Omega_n\}}(f - \mathbb{E}_{\mathrm{unif}\{\Omega_n\}} f)^2.$ Then we have the tightened Poincar\'e inequality 
 \[
 \frac{4}{n}\mathrm{Var}(f) \leq \langle f,Lf \rangle.
 \]
\end{theorem}
\begin{proof}
\[
\langle f,Lf \rangle = \sum_{B \in B_n} \frac{2|B|(n+1-|B|)}{n(n-1)}\widehat f(B)^2 \|\chi_B\|_2^2 = \sum_{\text{even } |B| \geq 2 }\frac{2|B|(n+1-|B|)}{n(n-1)}\widehat f(B)^2 \|\chi_B\|_2^2,
\]
since $\widehat{f}(B)=0$ for $B$  with $|B|$ odd by Theorem~\ref{eveness}. Moreover,
\[
\mathrm{Var}(f) = \langle f - \mathbb{E}_{\mathrm{unif}\{\Omega_n\}} f, f - \mathbb{E}_{\mathrm{unif}\{\Omega_n\}} f \rangle = \langle f - \widehat f(\emptyset) , f - \widehat f(\emptyset) \rangle = \sum_{\text{even } |B| \geq 2 }\widehat f(B)^2 \|\chi_B\|_2^2.
\]
By \cite[Lemma 2]{Filmus16}, $|B| \leq \frac{n}{2}$. Hence we have $\frac{2|B|(n+1-|B|)}{n(n-1)} \geq \frac{4}{n}$ for even $|B| \geq 2$ by direct calculation, completing the proof. 
\end{proof}


Another fact used in the Kwapie\'n-Lata\l a-Oleszkiewicz's proof of the Szarek's inequality is a pointwise bound for $Lf$. We now derive an analogue on the middle slice. 
\begin{lemma}\label{lemma-Lf}
Let $n\in \mathbb{N}$ be even and $n \geq 4$. Let $f$ be the function discussed above. Then
\[
Lf(x) \leq \frac{2}{n-1} f(x), \, \quad \text{for all } x \in \Omega_n.
\]
\end{lemma}
\begin{proof}
Denote by $S_n(x)=\sum_{i=1}^n a_ix_i$. Note that for $x \in \Omega_n$,
\begin{align*}
    \sum_{i<j} S_n^{(i,j)}(x) &= \sum_{i<j} \big(S_n(x) + (a_i-a_j)(x_j-x_i)\bigr) \\
    & = \binom{n}{2} S_n(x) + \frac{1}{2}\sum_{i\neq j}(a_ix_j -a_ix_i-a_jx_j+a_jx_i) \\
    & = \binom{n}{2} S_n(x) + \frac{1}{2}\sum_{i\neq j}a_ix_j - \frac{1}{2}\sum_{i\neq j}a_ix_i- \frac{1}{2}\sum_{i\neq j}a_jx_j+ \frac{1}{2}\sum_{i\neq j}a_jx_i\\
    & =  \binom{n}{2} S_n(x) + \frac{1}{2}\sum_{i=1}^n (a_i(\sum_{j=1}^n x_j)-a_ix_i) - \frac{1}{2}(n-1)S_n(x) \\ &- \frac{1}{2}(n-1)S_n(x) + \frac{1}{2}\sum_{j=1}^n (a_j(\sum_{i=1}^n x_i)-a_jx_j) \\
    & = \binom{n}{2} S_n(x) - \frac{1}{2}S_n(x) - \frac{1}{2}(n-1)S_n(x)- \frac{1}{2}(n-1)S_n(x) -  \frac{1}{2}S_n(x), \quad\text{since}\sum_{i=1}^n x_i=0 \\
    & = (\binom{n}{2} -n )S_n(x).
\end{align*}
Hence, for $x \in \Omega_n$
\begin{align*}
    Lf(x) &= |S_n(x)| - \frac{1}{\binom{n}{2}}\sum_{i<j}|S_n^{(i,j)}(x)| \\
    & \leq |S_n(x)| - \frac{1}{\binom{n}{2}}|\sum_{i<j}S_n^{(i,j)}(x)| \\
    & = |S_n(x)| - \frac{1}{\binom{n}{2}}|(\binom{n}{2}-n)S_n(x)| \\
    & = \frac{2}{n-1}f(x).
\end{align*}
\end{proof}

\begin{proof}[Proof of Theorem~\ref{Szarek}]
The upper bound follows from  Jensen's inequality. For the lower bound, we proceed as follows. Let $f(x)$ be the unique multilinear harmonic polynomial over $\mathbb{R}^n$ such that $f(x)=|\sum_{i=1}^n a_ix_i|$ for all $x\in \Omega_n$. Then by Theorem~\ref{Poincare} and Lemma~\ref{lemma-Lf}, we have
\[
\frac{4}{n}\mathrm{Var}(f) \leq \langle f,Lf\rangle \leq  \frac{2}{n-1} \|f\|_2^2.
\]
Since $\mathrm{Var}(f) = \|f\|_2^2 - (\mathbb{E}f)^2$, we obtain
\[
\mathbb{E}f \geq \sqrt{\frac{n-2}{2(n-1)}} \|f\|_2,
\]
completing the proof. 
\end{proof}
\begin{remark}
It would be interesting to explore further which symmetry condition of the distribution of $(x_1,\ldots,x_n)$ leads to Khintchine-type inequalities, without assuming that the $x_i$ are independent.
\end{remark}



\subsection{Proof of Theorem~\ref{sidon}}

The lower bound \(c\ge 1.28\) was proved in \cite{CloningerSteinerberger17}. The key estimate, established in \cite[Lemma~3]{CloningerSteinerberger17}, is that for all \(m,n\in\mathbb N\),
\[
c\ge b_{n,m}-\frac{2}{m}-\frac{1}{m^2},
\]
where \(b_{n,m}\) is defined in terms of a discrete set \(B_{n,m}\); see \cite[Lemma~3]{CloningerSteinerberger17} for the definition of $B_{n,m}$ and further details. Using a large-scale computation, Cloninger and Steinerberger \cite[Section 3.3]{CloningerSteinerberger17} showed that for \(n=24\) and \(m=50\),
\[
b_{n,m}\ge 1.28+\frac{2}{m}+\frac{1}{m^2},
\]
which yields \(c\ge 1.28\). Their computation required about \(20000\) CPU hours on a high-performance server.

Our improvement comes from a small refinement of the error term. Namely, we replace the estimate above by
\[
c\ge b_{n,m}-\frac{2}{m}-\frac{1}{2m^2}.
\]
Using the same computed value of \(b_{24,50}\), this gives
\[
c\ge 1.2802.
\]

The argument is very simple: the only change is in the bound for \(\varepsilon\ast\varepsilon\). In \cite[Lemma~3]{CloningerSteinerberger17}, one uses
\[
\|\varepsilon\ast\varepsilon\|_{L^\infty(\mathbb R)}\le \frac{1}{m^2}.
\]
Since \(\varepsilon\) is supported on \(\left[-\tfrac14,\tfrac14\right]\), which has length \(\tfrac12\), we have
\[
\|\varepsilon\|_{L^1(\mathbb R)}
\le \frac12\,\|\varepsilon\|_{L^\infty(\mathbb R)}
\le \frac{1}{2m}.
\]
Therefore,
\[
\|\varepsilon\ast\varepsilon\|_{L^\infty(\mathbb R)}
\le \|\varepsilon\|_{L^\infty(\mathbb R)}\|\varepsilon\|_{L^1(\mathbb R)}
\le \frac1m\cdot\frac{1}{2m}
= \frac{1}{2m^2}.
\]
Substituting this into the proof of \cite[Lemma~3]{CloningerSteinerberger17} yields
\[
c\ge b_{n,m}-\frac{2}{m}-\frac{1}{2m^2},
\]
as claimed.

\section*{Acknowledgments}
 P.I. acknowledges partial support from the NSF CAREER grant DMS-2152401, a Simons Fellowship, and a Humboldt Research Fellowship for Experienced Researchers. 

\appendix
\section{Chat Links}\label{appendix:A}

\begin{itemize}
    \item Theorem~\ref{Gaussian}. Chat link with Grok 4.20(Beta). \\
\url{https://grok.com/share/c2hhcmQtNA_452b1841-e7b1-4695-a6aa-cf3210365900?rid=f9834ecf-f716-4a48-81c6-15aefa0dcba7} \\

    \item Upper bound in Theorem~\ref{lem:sqrt3}. Chat link with Grok 4 Expert. \\
    \url{https://grok.com/share/c2hhcmQtNA_84e20fd0-d81a-4809-9728-83d37e88cd5f?rid=f6483ca7-f772-4e51-a8e2-f62926b6dfab}\\

    \item Lower bound in Theorem~\ref{lem:sqrt3}. Chat link with Grok 4 Heavy. \\
    \url{https://grok.com/share/c2hhcmQtNA_826731f3-afbc-42a5-aa61-acca3ec1324a?rid=de4a202d-357d-4ece-9954-16c0dc5951d5}\\

    \item Theorem~\ref{Szarek}. Chat link with Grok 4.30(Beta).\\
    \url{https://grok.com/share/c2hhcmQtNA_de1dfee6-b59f-4fa8-a439-aa176d2308c6}\\

    \item Theorem~\ref{sidon}. Chat link with Grok 4 Heavy.\\
    \url{https://grok.com/share/c2hhcmQtNQ_f4d17f80-4582-4679-b931-06277fd4cfd4?rid=a60436ae-eaba-4638-a0fd-47b231f19cd0}

\end{itemize}



 







\bibliographystyle{abbrv}
\begin{thebibliography}{10}


\bibitem{Ball1993}
K.~Ball.
\newblock The Reverse Isoperimetric Problem for Gaussian Measure.
\newblock \emph{Discrete \& Computational Geometry}, 10:411--420, 1993.

\bibitem{Bentkus03}
V.~Bentkus.
\newblock On the dependence of the Berry--Esseen bound on dimension.
\newblock {\em J. Statist. Plann. Inference}, Volume 113,
2003, pp.~385--402.


\bibitem{BubeckCoesterEldanGowersLeeLupsascaSawhneyScherrerSellkeSpearsUnutmazWeilYinZhivotovskiy25}
S.~Bubeck, C.~Coester, R.~Eldan, T.~Gowers, Y.~T.~Lee, A.~Lupsasca,
M.~Sawhney, R.~Scherrer, M.~Sellke, B.~K.~Spears, D.~Unutmaz,
K.~Weil, S.~Yin, N.~Zhivotovskiy.
\newblock Early science acceleration experiments with GPT-5.
\newblock {\em arXiv preprint arXiv:2511.16072},
2025.

\bibitem{BubeckDingEldanRacz16}
S.~Bubeck, J.~Ding, R.~Eldan, M.~Z.~R\'acz.
\newblock Testing for high-dimensional geometry in random graphs.
\newblock {\em Random Structures \& Algorithms}, Volume 49, Issue 3,
2016, pp.~503--532.

\bibitem{BurdenFairesBurden2016}
R.~L. Burden, J.~D. Faires, and A.~M. Burden,
\emph{Numerical Analysis},
10th ed.,
Cengage Learning, Boston, MA, 2016.

\bibitem{OptConstC1a}
$C_{1a}$ Sidon set autocorrelation constant. 
\newblock Optimization Constants in Mathematics.
\newblock \texttt{https://teorth.github.io/optimizationproblems/constants/1a.html},
accessed April 16, 2026.

\bibitem{CloningerSteinerberger17}
A.~Cloninger, S.~Steinerberger.
\newblock On suprema of autoconvolutions with an application to Sidon sets.
\newblock {\em Proceedings of the American Mathematical Society}, Volume 145, Issue 8,
2017, pp.~3191--3200.

\bibitem{EskenazisIvanisvili}
A.~Eskenazis and P.~Ivanisvili,
\newblock Polynomial inequalities on the Hamming cube,
\newblock {\em Probability Theory and Related Fields} {\bf 178} (2020), 235--287.

\bibitem{FengTrinhBinghamKangZhangKimBarretoSchildkrautJungSeoPaganoChervonyiHwangHouGukovTsaiChoiJinLiWuShiuShihLeLuong26}
T.~Feng, T.~Trinh, G.~Bingham, J.~Kang, S.~Zhang, S.-h.~Kim, K.~Barreto,
C.~Schildkraut, J.~Jung, J.~Seo, C.~Pagano, Y.~Chervonyi, D.~Hwang,
K.~Hou, S.~Gukov, C.-C.~Tsai, H.~Choi, Y.~Jin, W.-Y.~Li, H.-A.~Wu,
R.-A.~Shiu, Y.-S.~Shih, Q.~V.~Le, T.~Luong.
\newblock Semi-Autonomous Mathematics Discovery with Gemini: A Case Study on the Erd\H{o}s Problems.
\newblock {\em arXiv preprint arXiv:2601.22401},
2026.

\bibitem{Filmus16}
Y.~Filmus.
\newblock An Orthogonal Basis for Functions over a Slice of the Boolean Hypercube.
\newblock {\em The Electronic Journal of Combinatorics}, Volume 23, Issue 1, 2016.

\bibitem{SimonsOpen}
Y.~Filmus, H.~Hatami, S.~Heilman, E.~Mossel, R.~O'Donnell, S.~Sachdeva, A.~Wan, K.~Wimmer.
\newblock Real Analysis in Computer Science: A collection of Open Problems.
\newblock Simons Institute (2014). \url{https://simons.berkeley.edu/sites/default/files/openprobsmerged.pdf}.

\bibitem{Georgiev25}
B.~Georgiev, J.~G\'omez-Serrano, T.~Tao, A.~Z.~Wagner.
\newblock Mathematical exploration and discovery at scale.
\newblock {\em arXiv preprint arXiv:2511.02864},
2025.

\bibitem{Haagerup81}
U.~Haagerup.
\newblock The best constants in the Khintchine inequality.
\newblock {\em Studia Mathematica}, Volume 70,
1981, pp.~231--283.

\bibitem{HerscoviciSpektor20}
O.~Herscovici, S.~Spektor.
\newblock The best constant in the Khinchine inequality for slightly dependent random variables.
\newblock {\em arXiv preprint arXiv:1806.03562},
2020.

\bibitem{IvanisviliTkocz}
P.~Ivanisvili and T.~Tkocz,
\newblock Comparison of moments of Rademacher chaoses,
\newblock {\em Ark. Mat.} \textbf{57} (2019), no.~1, 121--128.

\bibitem{IvanisviliXie25}
P.~Ivanisvili, X.~Xie.
\newblock Counterexample to majority optimality in NICD with erasures.
\newblock {\em arXiv preprint arXiv:2510.20013},
2025.

\bibitem{JuGaoJiangWuSunChenYWangYWangZWangHePWuXiaoLiuDaiDong26}
H.~Ju, G.~Gao, J.~Jiang, B.~Wu, Z.~Sun, L.~Chen, Y.~Wang, Y.~Wang,
Z.~Wang, W.~He, P.~Wu, L.~Xiao, R.~Liu, B.~Dai, B.~Dong.
\newblock Automated Conjecture Resolution with Formal Verification.
\newblock {\em arXiv preprint arXiv:2604.03789},
2026.

\bibitem{KOS2008}
A.~Klivans, R.~O'Donnell, and R.~A.~Servedio.
\newblock Learning geometric concepts via Gaussian surface area.
\newblock In \emph{Proc.\ 49th IEEE Symposium on Foundations of Computer Science (FOCS)}, pages 541--550, 2008.

\bibitem{KwapienLatalaOleszkiewicz96}
S.~Kwapie\'n, R.~Lata{\l}a, K.~Oleszkiewicz.
\newblock Comparison of moments of sums of independent random variables and differential inequalities.
\newblock {\em J. Funct. Anal.}, Volume 136, Issue 1,
1996, pp.~258--268.

\bibitem{LarssonCohn}
L.~Larsson-Cohn,
\newblock \(L^p\)-norms of Hermite polynomials and an extremal problem on Wiener chaos,
\newblock {\em Ark. Mat.} \textbf{40} (2002), 133--144.

\bibitem{LatalaOleszkiewicz94}
R.~Lata{\l}a, K.~Oleszkiewicz.
\newblock On the best constant in the Khinchin-Kahane inequality.
\newblock {\em Studia Math.}, Volume 109, Issue 1,
1994, pp.~101--104.

\bibitem{MartinOBryant07}
G.~Martin, K.~O'Bryant.
\newblock The symmetric subset problem in continuous Ramsey theory.
\newblock {\em Experimental Mathematics}, Volume 16, Issue 2,
2007, pp.~145--165.

\bibitem{MartinOBryant09}
G.~Martin, K.~O'Bryant.
\newblock The supremum of autoconvolutions, with applications to additive number theory.
\newblock {\em Illinois Journal of Mathematics}, Volume 53, Issue 1,
2009, pp.~219--235.

\bibitem{MatolcsiVinuesa10}
M.~Matolcsi, C.~Vinuesa.
\newblock Improved bounds on the supremum of autoconvolutions.
\newblock {\em Journal of Mathematical Analysis and Applications}, Volume 372, Issue 2,
2010, pp.~439--447.

\bibitem{NadimpalliPascale2025}
S.~Nadimpalli and C.~Pascale.
\newblock On the Maximal Gaussian Perimeter of Convex Sets, Revisited.
\newblock Preprint, 2025. \newblock \href{https://arxiv.org/abs/2508.20079}{arXiv:2508.20079}.

\bibitem{NT23}
P.~Nayar and T.~Tkocz.
\newblock Extremal sections and projections of certain convex bodies: a survey.
\newblock {\em Harmonic Analysis and Convexity}, 343-390, 2023.

\bibitem{Nazarov2003}
F.~Nazarov.
\newblock On the maximal perimeter of a convex set in $\mathbb{R}^n$ with respect to a Gaussian measure.
\newblock In \emph{Geometric Aspects of Functional Analysis (2001--2002)}, pages 169--187.
\newblock Lecture Notes in Mathematics, Vol.\ 1807, Springer, 2003.

\bibitem{ODonnell}
R.~O'Donnell,
\newblock {\em Analysis of Boolean Functions},
\newblock Cambridge University Press, Cambridge, 2014.

\bibitem{Oleszkiewicz99}
K.~Oleszkiewicz.
\newblock Comparison of moments via Poincar\'e-type inequality.
\newblock In {\em Advances in stochastic inequalities} (Atlanta, GA, 1997),
pp.~135--148, {\em Contemp. Math.}, Volume 234,
Amer. Math. Soc., Providence, RI, 1999.

\bibitem{Raic2019}
M.~Rai\v{c}.
\newblock A multivariate Berry--Esseen theorem with explicit constants.
\newblock \emph{Bernoulli}, 25(4A):2824--2853, 2019.

\bibitem{SchinzelSchmidt02}
A.~Schinzel, W.~M.~Schmidt.
\newblock Comparison of \(L^1\)- and \(L^\infty\)-norms of squares of polynomials.
\newblock {\em Acta Arithmetica}, Volume 104, Issue 3,
2002, pp.~283--296.

\bibitem{Spektor16}
S.~Spektor.
\newblock Restricted Khinchine inequality.
\newblock {\em Canadian Mathematical Bulletin}, Volume 59,
2016, pp.~204--210.

\bibitem{Szarek76}
S.~J.~Szarek.
\newblock On the best constants in the Khinchin inequality.
\newblock {\em Studia Mathematica}, Volume 58,
1976, pp.~197--208.

\bibitem{Tao26}
T.~Tao.
\newblock A crowdsourced repository for optimization constants?
\newblock {\em What's new},
22 January 2026.
\newblock \url{https://terrytao.wordpress.com/2026/01/22/a-crowdsourced-repository-for-optimization-constants/}.

\bibitem{Tao25Mathstodon}
T.~Tao.
\newblock I encountered no issues with hallucinations or other AI-generated nonsense.
\newblock {\em Mathstodon post},
2 October 2025.
\newblock Status 115306550084415351.
\newblock \url{https://mathstodon.xyz/@tao/115306550084415351},
accessed April 20, 2026.

\bibitem{Wang25}
Y.~Wang, S.-R.~Su, Z.~Zeng, E.~Xu, L.~Ren, X.~Yang, Z.~Huang, X.~He,
L.~Ma, B.~Peng, H.~Cheng, P.~He, W.~Chen, S.~Wang, S.~S.~Du, Y.~Shen.
\newblock ThetaEvolve: Test-time Learning on Open Problems.
\newblock {\em arXiv preprint arXiv:2511.23473},
2025.

\bibitem{Wendel1948}
J.~G. Wendel,
\newblock Note on the gamma function,
\newblock \emph{Amer. Math. Monthly} \textbf{55} (1948), no.~9, 563--564.
\newblock doi:10.2307/2304460.

\bibitem{MO184286}
\newblock What is the minimal \(C_k\), such that every \(f\colon\{-1,1\}^n\to\mathbb{R}\) of degree at most \(k\) satisfies \(\|f\|_2\le C_k\|f\|_1\)?
\newblock {\em MathOverflow}, Question 184286, 2014.
\newblock Available at \url{https://mathoverflow.net/questions/184286/} (accessed April 3, 2026).

\bibitem{WoodruffCohenAddadJainMaoZuoBateniBranzeiBrennerChenFengFortnowFuGuanHadizadehHajiaghayiJafariRavizJavanmardKarthikKawarabayashiKumarLattanziLeeLiPanageasPaparasPrzybockiSubercaseauxSvenssonTaherijamWuYogevZadimoghaddamZhouMatiasManyikaMirrokni26}
D.~P.~Woodruff, V.~Cohen-Addad, L.~Jain, J.~Mao, S.~Zuo, M.~Bateni,
S.~Branzei, M.~P.~Brenner, L.~Chen, Y.~Feng, L.~Fortnow, G.~Fu, Z.~Guan,
Z.~Hadizadeh, M.~T.~Hajiaghayi, M.~JafariRaviz, A.~Javanmard, K.~C.~S.,
K.-i.~Kawarabayashi, R.~Kumar, S.~Lattanzi, E.~Lee, Y.~Li, I.~Panageas,
D.~Paparas, B.~Przybocki, B.~Subercaseaux, O.~Svensson, S.~Taherijam,
X.~Wu, E.~Yogev, M.~Zadimoghaddam, S.~Zhou, Y.~Matias, J.~Manyika,
V.~Mirrokni.
\newblock Accelerating Scientific Research with Gemini: Case Studies and Common Techniques.
\newblock {\em arXiv preprint arXiv:2602.03837},
2026.

\bibitem{Yuksekgonul26}
M.~Yuksekgonul, D.~Koceja, X.~Li, F.~Bianchi, J.~McCaleb, X.~Wang,
J.~Kautz, Y.~Choi, J.~Zou, C.~Guestrin, Y.~Sun.
\newblock Learning to Discover at Test Time.
\newblock {\em arXiv preprint arXiv:2601.16175},
2026.
\end{thebibliography}

\end{document}
%These results, following our previous work \cite{IvanisviliXie25}, also serve as part of our continuing effort to document the performance of frontier reasoning models and to investigate how best to exploit their potential in mathematical research. One advantage of our approach is that the questions we ask are natural mathematical questions posed to publicly available models, which is closer to the way mathematicians conduct research. It is also an approach that may invite broader public participation, since many people have small puzzles that can be formulated in mathematical language and may benefit from hints obtained by chatting with a reasoning model. Recently, a tremendous number of results have been obtained in a similar fashion; to name just a few \cite{IvanisviliXie25,BubeckCoesterEldanGowersLeeLupsascaSawhneyScherrerSellkeSpearsUnutmazWeilYinZhivotovskiy25,FengTrinhBinghamKangZhangKimBarretoSchildkrautJungSeoPaganoChervonyiHwangHouGukovTsaiChoiJinLiWuShiuShihLeLuong26, WoodruffCohenAddadJainMaoZuoBateniBranzeiBrennerChenFengFortnowFuGuanHadizadehHajiaghayiJafariRavizJavanmardKarthikKawarabayashiKumarLattanziLeeLiPanageasPaparasPrzybockiSubercaseauxSvenssonTaherijamWuYogevZadimoghaddamZhouMatiasManyikaMirrokni26}. 

If these models keep making significant progress, finding examples that are genuinely easy for humans to verify, but that these models fail to solve, may become very valuable yet challenging.
Core idea: our approach is speak to the machine in math language and get the new results also in math langue. 
Back ground: \\
- Previous systematic study is based on the coding system, whereas our approach only need to ask several nature questions. \\
- Most of reports are for the usage of  "brainstorm", writing codes. New results obtained through direction computation are rarely reported, and few reported case fails into a a limtied of fields combinatroics and learning theory. \\
- Our study report the direct approach, and spans fields in probability theory, convex geometry, classical analysis, additive combinatorics, showing what's more the current frontier models that can already do. \\
- Reporting failure case are rare. But we found it is important in understand the current frontier models and encourage the community to do so as well, as it maybe very important to figure out what theses system can not do in the future. \\

There have been many successfully mathematical discovery with large language models, most of them are reported for the usage of "brainstorm", writing codes, or integrated as a component designed specifically for mathematical discovery. 

Our notes focus on the report case for new results that are obtained directly from interaction from LLMs, i.e., 

On the other hand, new mathematical results obtained directly through LLMs were usually reported in Erdos problems. 

Here we select five instances from probability theory, convex geometry, classical analysis over different simple mathematics spaces, i.e., the hamming cube, $\mathbb{R}^n$, first $n$ positive integers. 

Through theses instances, we highlight that these models seems already very good at mathematical expolration within the landscape of current mathematical knowlege. 


When and why large language models become so effective? Because of the use of tools is internally cooperated with the reasoning process. 


%collecting such cautionary example becomes more and more important (mathematical mistkaes ai still make)

%Such errors make human verification necessary yet challenging in practice. Taken together, our successful and cautionary examples illustrate both the promise of frontier reasoning models in mathematical research and the urgent need for trustworthy, user-friendly verification tools.
%Large language models are often considered to make simple mistakes in mathematics. However, among recent frontier reasoning models, such mistakes may no longer be common. Through direct interaction with these systems, one may even discover new mathematical results. To illustrate this, we report several instances obtained through this approach and verified by the authors: 1.  an improved autocorrelation inequality and asymptotic size estimate for g-Sidon sets, 2. improved lower bound on the maximal Gaussian perimeter over convex bodies, 3. optimal balanced Szarek's inequality, 4. improved $L_2-L_1$ moment comparison inequality, 5. improved lower bound for the constant in the $L_2-L_1$ moment comparison inequality. 

%In contrast, despite their remarkable creativity and accuracy, these systems still make mistakes, sometimes in very subtle ways, making human involvement at the current stage both necessary and challenging. We report several such instances. Taken together, the successful and cautionary examples reveal both the promise of these systems and the practical difficulty of verification, underscoring the need for trustworthy, user-friendly verification tools in AI-assisted mathematical research.


%a rapidly growing demand for user-friendly verification systems in mathematical research.

% Through these instances, we try to outline how frontier reasoning models can help with tasks considered difficult, thereby strengthening our understanding for problems of current research interest. 



%Languages are fundamental way for humans to communicate, collaborate, and understand the world around us. Mathematical language is one of them that can describe almost everything in the world with arbitrary accuracy. Despite its benefits, modern mathematical research, due to its abstraction and complexity, is often limited to experts, forming a barrier to the dissemination of mathematical knowledge and its exposition in broad applications. With the rapid development of large language models (LLMs), there is growing hope that this barrier can be lowered by learning and advancing mathematical research through direct interaction with LLMs. To shed light on this potential, we present several instances in which direct human interactions with LLMs have contributed to advances in mathematical research across different fields and in a variety of ways: 1. an optimal balanced $L_1$ Khintchine inequality (balanced Szarek's inequality); 2. a first improved lower bound on the maximal Gaussian perimeter of convex sets  since Nazarov's 2003 result;  3. an improved lower bound of an autocorrelation inequality; 4....
%Theses instances are obtained from simply chatting with LLMs and the correctness is verified by the authors. The use of LLMs in theses instances spans mathematical exposition, refining hard analytical arguments, developing novel ideas suggested by the authors, identifying human oversights, and detecting gaps in probability, analysis and convex geometry. We also address that (1) asking the right questions is crucial from our experiences, and (2) the lack of correctness measure of LLM-generated results has been the current bottleneck in fully exploiting its mathematical ability. 
%large language models or reasoning generative models? the later word is closer to the mechanism of theses models. 
